\documentclass[12pt]{article}
\usepackage{amsmath,amssymb,amsthm}
\usepackage[margin=1in]{geometry}

\newtheorem{theorem}{Theorem}
\newtheorem{lemma}[theorem]{Lemma}

\title{Problem 228: Minkowski Sums}
\author{Project Euler Solution}
\date{}

\begin{document}
\maketitle

\section{Problem Statement}
Let $S_n$ be a regular $n$-sided polygon. Find the number of sides of
\[
T = S_2 + S_3 + S_4 + \cdots + S_{1864},
\]
where $+$ denotes the Minkowski sum.

\section{Mathematical Foundation}

\begin{theorem}[Minkowski Sum Edge Count]
The number of edges of the Minkowski sum of convex polygons equals the number of distinct edge normal directions across all summands.
\end{theorem}
\begin{proof}
The support function of the sum is $h_{P+Q}(\theta) = h_P(\theta) + h_Q(\theta)$. Edges correspond to maximal affine arcs of the support function. Two edges from different summands sharing a normal direction merge into one edge. Hence the edge count equals the count of distinct normals.
\end{proof}

\begin{lemma}[Normal Directions of $S_n$]
A regular $n$-gon has outward normals at angles $\theta_k = \pi(2k+1)/n$ for $k = 0, \ldots, n-1$.
\end{lemma}
\begin{proof}
Vertices at angles $2\pi k/n$; edge $k$ has midpoint direction $\pi(2k+1)/n$.
\end{proof}

\begin{theorem}[Reduced Fraction Characterization]
The direction $\pi p/q$ (with $\gcd(p,q) = 1$, $p$ odd, $0 < p < 2q$) appears in $S_n$ iff $q \mid n$ and $n/q$ is odd.
\end{theorem}
\begin{proof}
Write $\frac{2k+1}{n} = \frac{p}{q}$ in lowest terms. Since $2k+1$ is odd, $p$ is odd. Then $n = q(2k+1)/p$, requiring $q \mid n$ and $n/q$ odd.
\end{proof}

\begin{theorem}[Direction Count per Denominator]
For denominator $q$:
\begin{itemize}
\item $q$ odd: $\phi(q)$ directions.
\item $q$ even: $2\phi(q)$ directions.
\end{itemize}
\end{theorem}
\begin{proof}
Count odd $p \in (0, 2q)$ with $\gcd(p, q) = 1$. For odd $q$: $\gcd(p, 2q) = 1$ requires $p$ odd and $\gcd(p,q)=1$, giving $\phi(2q) = \phi(q)$ values. For even $q$: there are $2\phi(q)$ such odd values.
\end{proof}

\begin{theorem}[Existence]
Every direction with denominator $q \leq 1864$ is realized by some summand $S_n$ with $2 \leq n \leq 1864$.
\end{theorem}
\begin{proof}
For $q \geq 2$, $S_q$ is a summand. For $q = 1$, $S_3$ contains the direction.
\end{proof}

\section{Algorithm}
\begin{enumerate}
\item Compute $\phi(q)$ for $q = 1, \ldots, 1864$ via a sieve.
\item Sum: $\text{total} = \phi(1) + \sum_{q=2}^{1864} c(q)$, where $c(q) = \phi(q)$ if $q$ odd, $c(q) = 2\phi(q)$ if $q$ even.
\end{enumerate}

\section{Complexity Analysis}
\begin{itemize}
\item \textbf{Time:} $O(N \log\log N)$ for the totient sieve.
\item \textbf{Space:} $O(N)$.
\end{itemize}

\section{Answer}

\[
\boxed{86226}
\]

\end{document}
